% Generated by report/make_restated.py from theory/B1_preconditioned_saga.tex: do not edit (edit the note and rerun).
\begin{restated}{Remark}{rem:B1-dragomir}{B1}{0}{comparison with Bregman-SAGA}
Corollary~1 of \cite{dragomir2021fast} states, for Bregman-SAGA with a quadratic reference function
($\nabla^2h$ constant, which is exactly $P$-SAGA by Lemma~\ref{lem:B1-step}) and step $1/(8L_{f/h})$,
the contraction $\E[\psi_t]\le(1-\min\{\frac{1}{8\kappa_{f/h}},\frac1{2n}\})^t\psi_0$ of their potential
$\psi_t=\frac1{\eta}D_h(x^\star,x_t)+\frac n2\cdot\frac1n\sum_iD_{f_i}(\phi_i^t,x^\star)$, where
$L_{f/h}$ is the relative smoothness of every $f_i$ (their Assumption~3) and
$\kappa_{f/h}=L_{f/h}/\mu_{f/h}$ the relative condition number: these are our $R_P^2$ and $\kappaP$. The
rate has the same structure $1-\min\{c_1/n,\,c_2/\kappaP\}$ as \eqref{eq:B1-VP} (with $(c_1,c_2)=(\frac12,\frac18)$
instead of $(\frac13,\frac3{16})$, the difference coming from the step and from a potential expressed with
Bregman distances of the $f_i$ instead of squared norms of memory vectors).
\end{restated}
